% ══════════════════════════════════════════════════════════════════════════════
%  EXPERIMENT 3 — MEASURING TEMPERATURE AND TIME                       (Lesson 6)
% ══════════════════════════════════════════════════════════════════════════════
\experimentdiv{Measuring Temperature and Time}{Using a Bunsen burner safely \textbullet{} Reading a thermometer \textbullet{} Recording at regular intervals}{6}\label{sec:3}
\lessoncodes{VC2S8I01--I06 \textbullet{} VC2S8I08}

\vocabhead
\begin{ctab}
\begin{tabular}{|SW{4.2cm}|W{11.6cm}|}
  \hline\hd & \thead{Word} & \thead{Definition (copy from the slide)} \\ \hline
  & \scaf{\textbf{Temperature}} & \scaf{How hot or cold something is, measured in degrees Celsius (\textdegree C).} \\ \hline
  & \textbf{Thermometer} & \\ \hline
  & \textbf{Interval} & \\ \hline
  & \textbf{Safety flame} & \\ \hline
  & \textbf{Heating flame} & \\ \hline
\end{tabular}
\end{ctab}

\begin{lipanel}
\inquiryli{Use equipment safely to record temperature at regular time intervals in a table.}
\understandli{Describe how the temperature of water changes as it is heated.}
\sctitle
\begin{itemize}[leftmargin=1.4em, itemsep=0pt, topsep=1pt]
  \item I can light a Bunsen burner safely and switch between the safety flame and the heating flame.
  \item I can read a thermometer at eye level to the nearest degree.
  \item I can record the temperature of heating water every minute in a table, and describe the pattern.
\end{itemize}
\end{lipanel}

\begin{skillsfocus}
\focus{QP1}{Write a testable hypothesis: \emph{If the [IV] is \dots, the [DV] will \dots}}
\focus{PC2}{Name the independent, dependent and controlled variables.}
\focus{PC3}{Complete a risk assessment before you light the Bunsen burner.}
\focus{PC4}{Read the thermometer at eye level, to the nearest degree.}
\focus{PC6}{Take a reading at every whole minute, with its unit.}
\focus{PA5}{Describe the pattern: how temperature changes with time.}
\focus{EV2}{Compare your boiling point with the accepted value, 100\,\textdegree C.}
\end{skillsfocus}

% ─── TECHNIQUE ────────────────────────────────────────────────────────────────
\exppart{Technique spotlight}{Blue}

\begin{technique}{Using a Bunsen burner}{PC3}
\begin{minipage}[c]{0.62\linewidth}
\textbf{Before you light it}
\begin{itemize}[leftmargin=1.2em, itemsep=0pt, topsep=1pt]
  \item Safety glasses on, hair tied back, sleeves tucked in, bag under the bench.
  \item Stand the burner on a heatproof mat, away from paper and the bench edge.
  \item Check the gas tube for cracks, and push it firmly onto the gas tap.
\end{itemize}
\smallskip
\textbf{Lighting it}
\begin{enumerate}[leftmargin=1.6em, itemsep=0pt, topsep=1pt]
  \item Turn the collar so the air hole is \textbf{closed}.
  \item Light the match or striker \textbf{first}, and hold it beside the top of the barrel.
  \item Turn the gas tap on. You get the yellow \textbf{safety flame}.
  \item To heat, open the air hole to get the blue \textbf{heating flame}.
  \item Whenever you are not heating, close the air hole to show the safety flame.
  \item When you finish, turn the gas tap off.
\end{enumerate}
\end{minipage}\hfill
\begin{minipage}[c]{0.34\linewidth}\centering
\bunsenflame{safety}\hspace{0.5cm}\bunsenflame{heating}\\[3pt]
{\footnotesize\textbf{Safety flame}\hspace{0.6cm}\textbf{Heating flame}}\\[2pt]
{\scriptsize air hole closed: yellow, wavy, easy to see\\[1pt] air hole open: blue, hot, hard to see}
\end{minipage}
\end{technique}

\begin{youdobox}[The parts of a Bunsen burner]
\task Name each labelled part of the Bunsen burner. Choose from: \emph{air hole, barrel, base, collar, gas tube.}\skr{PC7}

\noindent
\begin{minipage}[c]{0.36\linewidth}\centering
\begin{tikzpicture}[scale=1.25]
  \eqBunsen
  \draw[eqlab] (0.58,2.0) -- (-0.4,2.0) node[eqtext, anchor=east, font=\bfseries\small] {A};
  \draw[eqlab] (1.1,0.75) -- (1.9,0.95) node[eqtext, anchor=west, font=\bfseries\small] {B};
  \draw[eqlab] (0.8,0.52) -- (1.9,0.45) node[eqtext, anchor=west, font=\bfseries\small] {C};
  \draw[eqlab] (1.6,0.1) -- (1.9,-0.1) node[eqtext, anchor=west, font=\bfseries\small] {D};
  \draw[eqlab] (-0.4,0.11) -- (-0.6,-0.3) node[eqtext, anchor=north, font=\bfseries\small] {E};
\end{tikzpicture}
\end{minipage}\hfill
\begin{minipage}[c]{0.6\linewidth}
\begin{ctab}
\begin{tabular}{|SZ{1.2cm}|W{6.6cm}|}
  \hline\hd & \thead{Part} & \thead{Name} \\ \hline
  & \scaf{A} & \scaf{Barrel} \\ \hline
  & B & \\ \hline
  & C & \\ \hline
  & D & \\ \hline
  & E & \\ \hline
\end{tabular}
\end{ctab}
\end{minipage}
\end{youdobox}


\begin{technique}{Reading a thermometer}{PC4}
\begin{enumerate}[leftmargin=1.6em, itemsep=1pt, topsep=0pt]
  \item Keep the bulb in the middle of the liquid. It must not touch the bottom or side of the beaker, which are hotter than the water.
  \item Work out the scale: the difference between two labelled marks, divided by the number of spaces between them.
  \item Wait for the liquid in the thermometer to stop moving.
  \item Bend down so your eye is level with the top of the coloured liquid, then read the scale.
  \item Record the reading and the unit, \textdegree C.
\end{enumerate}
\end{technique}

\begin{workedex}{Reading a thermometer}
\begin{minipage}[c]{0.66\linewidth}
Read the temperature shown.

\smallskip
\textbf{Scale:} the marks 40 and 50 are 10\,\textdegree C apart, with 10 spaces between them. One small division is $10 \div 10 = 1$\,\textdegree C.

\smallskip
\textbf{Reading:} the top of the liquid is 7 small divisions above 40, so the temperature is $40 + 7 \times 1 = $ \textbf{47\,\textdegree C}.
\end{minipage}\hfill
\begin{minipage}[c]{0.3\linewidth}\centering
\readtherm{20}{60}{10}{10}{47}
\end{minipage}
\end{workedex}

\begin{wedo}{Reading a thermometer}
\begin{minipage}[c]{0.66\linewidth}
\pt{a}What is the value of one small division?
\workbox{2.2cm}
\pt{b}Read the temperature.
\answerbox{1cm}
\end{minipage}\hfill
\begin{minipage}[c]{0.3\linewidth}\centering
\readtherm{0}{40}{10}{5}{24}
\end{minipage}
\end{wedo}

\begin{youdobox}[Reading thermometers]
\task Read each thermometer. Give the value of one small division and the temperature, with units.\skr{PC4}

\begin{center}
\begin{tabular}{ccc}
  \readtherm{60}{110}{10}{5}{98} & \hspace{1.6cm}\readtherm{0}{50}{10}{10}{19} & \hspace{1.6cm}\readtherm{-10}{20}{10}{10}{-3} \\[4pt]
  \textbf{(a)} & \hspace{1.6cm}\textbf{(b)} & \hspace{1.6cm}\textbf{(c)} \\
\end{tabular}
\end{center}
\begin{ctab}
\begin{tabular}{|SZ{1.9cm}|Z{6.0cm}|Z{6.0cm}|}
  \hline\hd & \thead{Thermometer} & \thead{Value of one small division (\textdegree C)} & \thead{Temperature (\textdegree C)} \\ \hline
  & \scaf{e.g.} & \scaf{1} & \scaf{47} \\ \hline
  & (a) & & \\ \hline
  & (b) & & \\ \hline
  & (c) & & \\ \hline
\end{tabular}
\end{ctab}
\end{youdobox}

\begin{technique}{Timing at regular intervals}{PC8}
Share three jobs in your group. The \textbf{timer} starts the stopwatch the moment heating starts and calls ``Now!'' at every whole minute. The \textbf{reader} reads the thermometer on each call. The \textbf{recorder} writes the reading in the table straight away. Swap jobs halfway through. Even intervals let you compare one minute with the next.
\end{technique}

% ─── PLAN AND SAFETY ─────────────────────────────────────────────────────────
\exppart{Plan and safety}{Teal}

\textbf{Aim:} To investigate the effect of changing the heating time on the temperature of water.

\needspace{10\baselineskip}
\task Write a hypothesis for this experiment. Complete the sentence frame, then add a reason if you can.\skr{QP1}

\scaf{If the heating time is increased, the temperature of the water will \dots}
\answerbox{5cm}

\needspace{12\baselineskip}
\task Identify the variables in this experiment.\skr{PC2}
\begin{ctab}
\begin{tabular}{|SW{4.6cm}|W{4.6cm}|W{6.0cm}|}
  \hline\hd & \thead{Type of variable} & \thead{What it means} & \thead{In this experiment} \\ \hline
  & \scaf{Independent} & \scaf{What we change or choose} & \scaf{Time (minutes)} \\ \hline
  & Dependent & What we measure & \\ \hline
  & Controlled (1) & What we keep the same & \\ \hline
  & Controlled (2) & What we keep the same & \\ \hline
\end{tabular}
\end{ctab}

\needspace{14\baselineskip}
\task Complete the risk assessment before you light the Bunsen burner. For each risk, say how to prevent it and what to do if it happens. Your teacher must initial it before you start.\skr{PC3}\par\nobreak\hfill\textbf{Teacher's initials:}\ \underline{\hspace{3cm}}
\begin{ctab}
\begin{tabular}{|SW{3.8cm}|W{5.4cm}|W{5.8cm}|}
  \hline\hd & \thead{Risk (what could go wrong)} & \thead{Prevention} & \thead{Responding (if it happens)} \\ \hline
  & \scaf{Burns from hot water or glassware} & \scaf{Do not touch the beaker or tripod until cool; stand up while heating} & \scaf{Cool the burn under running water for 20 minutes; tell the teacher} \\ \hline
  & Hair or clothing catches the Bunsen flame & & \\ \hline
  & Gas escapes from the tap & & \\ \hline
  & The glass thermometer breaks & & \\ \hline
\end{tabular}
\end{ctab}

% ─── METHOD ──────────────────────────────────────────────────────────────────
\exppart{Method}{Teal}

\subsubsection*{Equipment}
\begin{multicols}{2}
\begin{itemize}[leftmargin=1.2em, itemsep=0pt, topsep=0pt]
  \item Bunsen burner and heatproof mat
  \item tripod and gauze mat
  \item 250\,mL beaker
  \item 100\,mL measuring cylinder
  \item thermometer (\textminus10 to 110\,\textdegree C)
  \item stopwatch
  \item matches or a striker
  \item safety glasses
\end{itemize}
\end{multicols}

\begin{center}
\begin{tikzpicture}[scale=0.95]
  \draw[eq] (-0.6,-0.12) rectangle (4.0,0);
  \begin{scope}[xshift=0.9cm]\eqBunsen\end{scope}
  \draw[Blue!35, line width=0.6pt, fill=Blue!12] (1.7,2.3) .. controls (1.4,2.7) and (1.55,3.1) .. (1.7,3.35) .. controls (1.85,3.1) and (2.0,2.7) .. (1.7,2.3);
  \draw[eq] (0.25,3.45) rectangle (3.15,3.57);
  \draw[eq] (0.55,3.45) -- (0.25,0) (2.85,3.45) -- (3.15,0);
  \draw[Sky!60, line width=0.6pt] (0.8,5.0) -- (2.6,5.0);
  \begin{scope}[xshift=0.8cm, yshift=3.57cm]\eqBeaker\end{scope}
  \draw[eq] (1.7,4.05) circle (0.15);
  \draw[eq] (1.62,4.19) -- (1.62,7.0) (1.78,4.19) -- (1.78,7.0);
  \draw[Ink, line width=1.3pt] (1.7,4.2) -- (1.7,5.6);
  \eqlab{(1.78,6.6)}{(4.6,6.6)}{thermometer}
  \eqlab{(2.6,5.8)}{(4.6,5.8)}{beaker}
  \eqlab{(2.2,5.0)}{(4.6,5.0)}{water}
  \eqlab{(3.15,3.51)}{(4.6,3.8)}{gauze mat}
  \eqlab{(3.05,1.5)}{(4.6,2.6)}{tripod}
  \eqlab{(1.92,1.4)}{(4.6,1.4)}{Bunsen burner}
  \eqlab{(4.0,-0.06)}{(4.6,0.3)}{heatproof mat}
  \eqlabl{(0.25,0.11)}{(-0.5,0.8)}{gas}
\end{tikzpicture}
\expfigcap{Scientific diagram of the equipment set up for heating water.}
\end{center}

\begin{enumerate}[itemsep=2pt]
  \item Put on safety glasses. Set up the heatproof mat, Bunsen burner and tripod with the gauze mat on top, as in the diagram.
  \item Measure 150\,mL of water with the measuring cylinder and pour it into the beaker. Stand the beaker in the middle of the gauze mat.
  \item Put the thermometer in the water. Read the starting temperature and record it at 0 minutes.
  \item Light the Bunsen burner (safety flame), slide it under the tripod, then open the air hole to the heating flame. Start the stopwatch.
  \item Read the temperature at every whole minute for 10 minutes. Record each reading and anything you observe (bubbles, steam, sounds).
  \item Close the air hole, turn off the gas, and leave everything to cool before you touch it.
\end{enumerate}

% ─── RESULTS ─────────────────────────────────────────────────────────────────
\exppart{Results}{Amber}

\subsubsection*{The temperature of water as it is heated \hfill \sk{PA1}}
\begin{ctab}
\begin{tabular}{|SZ{2.6cm}|Z{4.0cm}|W{8.8cm}|}
  \hline\hd & \thead{Time (min)} & \thead{Temperature (\textdegree C)} & \thead{Observations} \\ \hline
  & \scaf{e.g. 0} & \scaf{17} & \scaf{Water still and clear} \\ \hline
  & 0 & & \\ \hline
  & 1 & & \\ \hline
  & 2 & & \\ \hline
  & 3 & & \\ \hline
  & 4 & & \\ \hline
  & 5 & & \\ \hline
  & 6 & & \\ \hline
  & 7 & & \\ \hline
  & 8 & & \\ \hline
  & 9 & & \\ \hline
  & 10 & & \\ \hline
\end{tabular}
\end{ctab}

\begin{infobox}{Keep this data}
You will draw your first line graph from these results in The Particle Model. Make sure every reading has been recorded.\skr{PA2}
\end{infobox}

% ─── ANALYSIS ────────────────────────────────────────────────────────────────
\exppart{Analysis and evaluation}{Coral}

\task Describe the pattern in your results: how did the temperature of the water change as time went on? Quote at least two readings from your table.\skr{PA5}

\scaf{As the time increased, the temperature \dots. At 0 minutes it was \dots\,\textdegree C and at \dots\ minutes it was \dots\,\textdegree C.}
\answerbox{5cm}

\needspace{10\baselineskip}
\task Calculate how much the temperature rose in the first 3 minutes (0 to 3) and in the last 3 minutes (7 to 10). Was the rise the same?\skr{PA3}
\workbox{4cm}

\needspace{14\baselineskip}
\task Water boils at 100\,\textdegree C at sea level.\skr{EV2}
\pt{a}What was the highest temperature you recorded?
\answerbox{1cm}
\pt{b}Compare it with 100\,\textdegree C, and suggest one reason why your value might be different.
\answerbox{5cm}

\needspace{10\baselineskip}
\task Was there any reading that did not fit the pattern? If so, give it and suggest why it happened. If not, say how you can tell.\skr{PA7}
\answerbox{5cm}

\needspace{10\baselineskip}
\task Did your results support your hypothesis? Use your data as evidence. (Results \emph{support} a hypothesis; they never \emph{prove} it.)\skr{EV1}
\answerbox{5cm}

\needspace{10\baselineskip}
\task Why must you wait for the thermometer reading to stop moving, and read it at eye level?\skr{PC4}
\answerbox{5cm}

\needspace{10\baselineskip}
\task Suggest one change to the method that would make your results more precise, and explain how it would help.\skr{EV3}
\answerbox{5cm}

\begin{technique}{Writing a conclusion}{EV1}
A conclusion is two or three sentences in the third person:
\begin{enumerate}[leftmargin=1.6em, itemsep=1pt, topsep=3pt]
  \item \textbf{The general trend:} how the dependent variable changed as the independent variable changed.
  \item \textbf{Numbers that back it up:} the first and last values, or the biggest change.
  \item \textbf{The hypothesis:} say whether the results \emph{support} it. Never write that it was \emph{proven}.
\end{enumerate}
\emph{Example: As the amount of sunlight increased, the height of the plant increased, from 4\,cm to 11\,cm. This supports the hypothesis.}
\end{technique}

\needspace{12\baselineskip}
\task Write a conclusion for this experiment in the third person. State the general trend, back it with numbers from your table, and say whether your hypothesis was supported.\skr{CO1}

\scaf{As the heating time increased, the temperature of the water \dots, from \dots\,\textdegree C to \dots\,\textdegree C. This supports / does not support the hypothesis.}
\answerbox{5.5cm}

\begin{yr8box}
\begin{keepext}
\exttask Once water starts to boil, its temperature stops rising, even though the Bunsen burner is still heating it. Suggest where the energy from the flame is going. (You will find out in The Particle Model.)
\answerbox{5cm}
\end{keepext}
\end{yr8box}

% ─── SKILL CHECK ─────────────────────────────────────────────────────────────
\exppart{Skill check}{NavyBlue}

\begin{skillcheck}
\scrow{QP1}{I can write a hypothesis as \emph{If the [IV] is \dots, the [DV] will \dots}}
\scrow{PC2}{I can name the independent, dependent and controlled variables.}
\scrow{PC3}{I can complete a risk assessment, and light and use a Bunsen burner safely.}
\scrow{PC4}{I can read a thermometer at eye level, to the nearest degree.}
\scrow{PC6}{I recorded a reading at every whole minute, with its unit.}
\scrow{PC8}{I used a stopwatch to keep the intervals even.}
\scrow{PA5}{I can describe the pattern in my results, quoting data.}
\scrow{EV1}{I can say whether my results supported my hypothesis (supported, not proven).}
\scrow{EV2}{I can compare my result with the accepted value.}
\scrow{CO1}{I can write a conclusion in the third person.}
\end{skillcheck}

\task Fill in the \emph{End of unit} column of your skills tracker at the front of this booklet. Which skill has improved most since Experiment 1, and which one will you work on in The Particle Model?
\answerbox{5cm}

\end{document}
